% Options for packages loaded elsewhere
\PassOptionsToPackage{unicode}{hyperref}
\PassOptionsToPackage{hyphens}{url}
\documentclass[
]{article}
\usepackage{xcolor}
\usepackage[margin=1in]{geometry}
\usepackage{amsmath,amssymb}
\setcounter{secnumdepth}{-\maxdimen} % remove section numbering
\usepackage{iftex}
\ifPDFTeX
  \usepackage[T1]{fontenc}
  \usepackage[utf8]{inputenc}
  \usepackage{textcomp} % provide euro and other symbols
\else % if luatex or xetex
  \usepackage{unicode-math} % this also loads fontspec
  \defaultfontfeatures{Scale=MatchLowercase}
  \defaultfontfeatures[\rmfamily]{Ligatures=TeX,Scale=1}
\fi
\usepackage{lmodern}
\ifPDFTeX\else
  % xetex/luatex font selection
\fi
% Use upquote if available, for straight quotes in verbatim environments
\IfFileExists{upquote.sty}{\usepackage{upquote}}{}
\IfFileExists{microtype.sty}{% use microtype if available
  \usepackage[]{microtype}
  \UseMicrotypeSet[protrusion]{basicmath} % disable protrusion for tt fonts
}{}
\makeatletter
\@ifundefined{KOMAClassName}{% if non-KOMA class
  \IfFileExists{parskip.sty}{%
    \usepackage{parskip}
  }{% else
    \setlength{\parindent}{0pt}
    \setlength{\parskip}{6pt plus 2pt minus 1pt}}
}{% if KOMA class
  \KOMAoptions{parskip=half}}
\makeatother
\usepackage{color}
\usepackage{fancyvrb}
\newcommand{\VerbBar}{|}
\newcommand{\VERB}{\Verb[commandchars=\\\{\}]}
\DefineVerbatimEnvironment{Highlighting}{Verbatim}{commandchars=\\\{\}}
% Add ',fontsize=\small' for more characters per line
\usepackage{framed}
\definecolor{shadecolor}{RGB}{248,248,248}
\newenvironment{Shaded}{\begin{snugshade}}{\end{snugshade}}
\newcommand{\AlertTok}[1]{\textcolor[rgb]{0.94,0.16,0.16}{#1}}
\newcommand{\AnnotationTok}[1]{\textcolor[rgb]{0.56,0.35,0.01}{\textbf{\textit{#1}}}}
\newcommand{\AttributeTok}[1]{\textcolor[rgb]{0.13,0.29,0.53}{#1}}
\newcommand{\BaseNTok}[1]{\textcolor[rgb]{0.00,0.00,0.81}{#1}}
\newcommand{\BuiltInTok}[1]{#1}
\newcommand{\CharTok}[1]{\textcolor[rgb]{0.31,0.60,0.02}{#1}}
\newcommand{\CommentTok}[1]{\textcolor[rgb]{0.56,0.35,0.01}{\textit{#1}}}
\newcommand{\CommentVarTok}[1]{\textcolor[rgb]{0.56,0.35,0.01}{\textbf{\textit{#1}}}}
\newcommand{\ConstantTok}[1]{\textcolor[rgb]{0.56,0.35,0.01}{#1}}
\newcommand{\ControlFlowTok}[1]{\textcolor[rgb]{0.13,0.29,0.53}{\textbf{#1}}}
\newcommand{\DataTypeTok}[1]{\textcolor[rgb]{0.13,0.29,0.53}{#1}}
\newcommand{\DecValTok}[1]{\textcolor[rgb]{0.00,0.00,0.81}{#1}}
\newcommand{\DocumentationTok}[1]{\textcolor[rgb]{0.56,0.35,0.01}{\textbf{\textit{#1}}}}
\newcommand{\ErrorTok}[1]{\textcolor[rgb]{0.64,0.00,0.00}{\textbf{#1}}}
\newcommand{\ExtensionTok}[1]{#1}
\newcommand{\FloatTok}[1]{\textcolor[rgb]{0.00,0.00,0.81}{#1}}
\newcommand{\FunctionTok}[1]{\textcolor[rgb]{0.13,0.29,0.53}{\textbf{#1}}}
\newcommand{\ImportTok}[1]{#1}
\newcommand{\InformationTok}[1]{\textcolor[rgb]{0.56,0.35,0.01}{\textbf{\textit{#1}}}}
\newcommand{\KeywordTok}[1]{\textcolor[rgb]{0.13,0.29,0.53}{\textbf{#1}}}
\newcommand{\NormalTok}[1]{#1}
\newcommand{\OperatorTok}[1]{\textcolor[rgb]{0.81,0.36,0.00}{\textbf{#1}}}
\newcommand{\OtherTok}[1]{\textcolor[rgb]{0.56,0.35,0.01}{#1}}
\newcommand{\PreprocessorTok}[1]{\textcolor[rgb]{0.56,0.35,0.01}{\textit{#1}}}
\newcommand{\RegionMarkerTok}[1]{#1}
\newcommand{\SpecialCharTok}[1]{\textcolor[rgb]{0.81,0.36,0.00}{\textbf{#1}}}
\newcommand{\SpecialStringTok}[1]{\textcolor[rgb]{0.31,0.60,0.02}{#1}}
\newcommand{\StringTok}[1]{\textcolor[rgb]{0.31,0.60,0.02}{#1}}
\newcommand{\VariableTok}[1]{\textcolor[rgb]{0.00,0.00,0.00}{#1}}
\newcommand{\VerbatimStringTok}[1]{\textcolor[rgb]{0.31,0.60,0.02}{#1}}
\newcommand{\WarningTok}[1]{\textcolor[rgb]{0.56,0.35,0.01}{\textbf{\textit{#1}}}}
\usepackage{graphicx}
\makeatletter
\newsavebox\pandoc@box
\newcommand*\pandocbounded[1]{% scales image to fit in text height/width
  \sbox\pandoc@box{#1}%
  \Gscale@div\@tempa{\textheight}{\dimexpr\ht\pandoc@box+\dp\pandoc@box\relax}%
  \Gscale@div\@tempb{\linewidth}{\wd\pandoc@box}%
  \ifdim\@tempb\p@<\@tempa\p@\let\@tempa\@tempb\fi% select the smaller of both
  \ifdim\@tempa\p@<\p@\scalebox{\@tempa}{\usebox\pandoc@box}%
  \else\usebox{\pandoc@box}%
  \fi%
}
% Set default figure placement to htbp
\def\fps@figure{htbp}
\makeatother
\setlength{\emergencystretch}{3em} % prevent overfull lines
\providecommand{\tightlist}{%
  \setlength{\itemsep}{0pt}\setlength{\parskip}{0pt}}
\usepackage{bookmark}
\IfFileExists{xurl.sty}{\usepackage{xurl}}{} % add URL line breaks if available
\urlstyle{same}
\hypersetup{
  pdftitle={SIR and SEIR models with deSolve},
  hidelinks,
  pdfcreator={LaTeX via pandoc}}

\title{SIR and SEIR models with deSolve}
\author{}
\date{\vspace{-2.5em}}

\begin{document}
\maketitle

\textbf{Workflow lab:} Goal → Approach → Execution → Check → Report.

\subsection{Learning objectives}\label{learning-objectives}

\begin{itemize}
\tightlist
\item
  Write a compartmental SIR / SEIR model as a system of ODEs.
\item
  Solve it with \texttt{deSolve::ode} and plot the trajectories.
\item
  Compute the basic reproduction number R₀ from model parameters and
  illustrate its effect.
\end{itemize}

\subsection{Prerequisites}\label{prerequisites}

Basic calculus intuition and comfort with differential-equation
notation.

\subsection{Background}\label{background}

Compartmental models partition a population into states ---
\emph{Susceptible}, \emph{Exposed}, \emph{Infectious}, \emph{Recovered}
--- and describe flows between them with ordinary differential
equations. The SIR model has three compartments and two parameters
(transmission rate β, recovery rate γ); the SEIR model adds an
exposed-but-not-yet-infectious class and a progression rate σ. These
models are caricatures, but they capture the shape of an epidemic ---
exponential rise, peak, decline --- surprisingly well, and they make the
role of R₀ = β/γ visible.

The real power of compartmental models in practice is not prediction but
counterfactual reasoning: \emph{what if} we had vaccinated half the
population on day zero, or \emph{what if} the infectious period were a
day shorter? Tuning β and γ to reflect interventions is the bread and
butter of public-health modelling.

\subsection{Setup}\label{setup}

\begin{Shaded}
\begin{Highlighting}[]
\FunctionTok{library}\NormalTok{(tidyverse)}
\FunctionTok{library}\NormalTok{(deSolve)}
\FunctionTok{set.seed}\NormalTok{(}\DecValTok{42}\NormalTok{)}
\FunctionTok{theme\_set}\NormalTok{(}\FunctionTok{theme\_minimal}\NormalTok{(}\AttributeTok{base\_size =} \DecValTok{12}\NormalTok{))}
\end{Highlighting}
\end{Shaded}

\subsection{1. Goal}\label{goal}

Simulate an SIR outbreak in a population of 1 million, for R₀ values
1.5, 2.5, and 4.

\subsection{2. Approach}\label{approach}

\begin{Shaded}
\begin{Highlighting}[]
  \FunctionTok{with}\NormalTok{(}\FunctionTok{as.list}\NormalTok{(}\FunctionTok{c}\NormalTok{(state, params)), \{}
\NormalTok{    N  }\OtherTok{\textless{}{-}}\NormalTok{ S }\SpecialCharTok{+}\NormalTok{ I }\SpecialCharTok{+}\NormalTok{ R}
\NormalTok{    dS }\OtherTok{\textless{}{-}} \SpecialCharTok{{-}}\NormalTok{beta }\SpecialCharTok{*}\NormalTok{ S }\SpecialCharTok{*}\NormalTok{ I }\SpecialCharTok{/}\NormalTok{ N}
\NormalTok{    dI }\OtherTok{\textless{}{-}}\NormalTok{  beta }\SpecialCharTok{*}\NormalTok{ S }\SpecialCharTok{*}\NormalTok{ I }\SpecialCharTok{/}\NormalTok{ N }\SpecialCharTok{{-}}\NormalTok{ gamma }\SpecialCharTok{*}\NormalTok{ I}
\NormalTok{    dR }\OtherTok{\textless{}{-}}\NormalTok{  gamma }\SpecialCharTok{*}\NormalTok{ I}
    \FunctionTok{list}\NormalTok{(}\FunctionTok{c}\NormalTok{(dS, dI, dR))}
\NormalTok{  \})}
\ErrorTok{\}}
\end{Highlighting}
\end{Shaded}

\subsection{3. Execution}\label{execution}

\begin{Shaded}
\begin{Highlighting}[]
\NormalTok{  params }\OtherTok{\textless{}{-}} \FunctionTok{c}\NormalTok{(}\AttributeTok{beta =}\NormalTok{ R0 }\SpecialCharTok{*}\NormalTok{ gamma, }\AttributeTok{gamma =}\NormalTok{ gamma)}
\NormalTok{  state  }\OtherTok{\textless{}{-}} \FunctionTok{c}\NormalTok{(}\AttributeTok{S =}\NormalTok{ N }\SpecialCharTok{{-}}\NormalTok{ I0, }\AttributeTok{I =}\NormalTok{ I0, }\AttributeTok{R =} \DecValTok{0}\NormalTok{)}
\NormalTok{  times  }\OtherTok{\textless{}{-}} \FunctionTok{seq}\NormalTok{(}\DecValTok{0}\NormalTok{, days, }\AttributeTok{by =} \DecValTok{1}\NormalTok{)}
\NormalTok{  out }\OtherTok{\textless{}{-}} \FunctionTok{as.data.frame}\NormalTok{(}\FunctionTok{ode}\NormalTok{(state, times, sir, params))}
\NormalTok{  out}\SpecialCharTok{$}\NormalTok{R0 }\OtherTok{\textless{}{-}}\NormalTok{ R0}
\NormalTok{  out}
\ErrorTok{\}}

\NormalTok{trajectories }\OtherTok{\textless{}{-}} \FunctionTok{bind\_rows}\NormalTok{(}
  \FunctionTok{run\_sir}\NormalTok{(}\FloatTok{1.5}\NormalTok{), }\FunctionTok{run\_sir}\NormalTok{(}\FloatTok{2.5}\NormalTok{), }\FunctionTok{run\_sir}\NormalTok{(}\FloatTok{4.0}\NormalTok{)}
\NormalTok{) }\SpecialCharTok{|\textgreater{}}
  \FunctionTok{pivot\_longer}\NormalTok{(}\FunctionTok{c}\NormalTok{(S, I, R), }\AttributeTok{names\_to =} \StringTok{"compartment"}\NormalTok{, }\AttributeTok{values\_to =} \StringTok{"n"}\NormalTok{)}
\end{Highlighting}
\end{Shaded}

\begin{Shaded}
\begin{Highlighting}[]
\FunctionTok{ggplot}\NormalTok{(trajectories, }\FunctionTok{aes}\NormalTok{(time, n, }\AttributeTok{colour =}\NormalTok{ compartment)) }\SpecialCharTok{+}
  \FunctionTok{geom\_line}\NormalTok{(}\AttributeTok{linewidth =} \FloatTok{0.8}\NormalTok{) }\SpecialCharTok{+}
  \FunctionTok{facet\_wrap}\NormalTok{(}\SpecialCharTok{\textasciitilde{}}\NormalTok{ R0, }\AttributeTok{labeller =} \FunctionTok{labeller}\NormalTok{(}\AttributeTok{R0 =}\NormalTok{ \textbackslash{}(x) }\FunctionTok{paste0}\NormalTok{(}\StringTok{"R0 = "}\NormalTok{, x))) }\SpecialCharTok{+}
  \FunctionTok{labs}\NormalTok{(}\AttributeTok{x =} \StringTok{"Day"}\NormalTok{, }\AttributeTok{y =} \StringTok{"People"}\NormalTok{)}
\end{Highlighting}
\end{Shaded}

\subsection{4. Check}\label{check}

Peak prevalence and attack rate should rise with R₀.

\begin{Shaded}
\begin{Highlighting}[]
  \FunctionTok{filter}\NormalTok{(compartment }\SpecialCharTok{==} \StringTok{"I"}\NormalTok{) }\SpecialCharTok{|\textgreater{}}
  \FunctionTok{group\_by}\NormalTok{(R0) }\SpecialCharTok{|\textgreater{}}
  \FunctionTok{summarise}\NormalTok{(}\AttributeTok{peak\_I =} \FunctionTok{max}\NormalTok{(n), }\AttributeTok{peak\_day =}\NormalTok{ time[}\FunctionTok{which.max}\NormalTok{(n)]) }\SpecialCharTok{|\textgreater{}}
  \FunctionTok{arrange}\NormalTok{(R0)}
\end{Highlighting}
\end{Shaded}

\subsection{5. Report}\label{report}

\begin{quote}
An SIR compartmental model was simulated for a population of one million
with a 7-day mean infectious period across three basic reproduction
numbers (R₀ = 1.5, 2.5, 4). Peak prevalence and the final attack rate
increased sharply with R₀, illustrating how small changes in
transmission translate into large changes in epidemic size. Calibrated
analogues of this model underpin much real-world outbreak forecasting.
\end{quote}

\subsection{Common pitfalls}\label{common-pitfalls}

\begin{itemize}
\tightlist
\item
  Interpreting SIR trajectories as forecasts rather than scenarios.
\item
  Forgetting that a constant β across time ignores interventions.
\item
  Comparing peak prevalence across parameter settings without holding N
  and I₀ constant.
\end{itemize}

\subsection{Further reading}\label{further-reading}

\begin{itemize}
\tightlist
\item
  Keeling MJ, Rohani P. \emph{Modeling Infectious Diseases in Humans and
  Animals.}
\item
  Soetaert K et al.~\emph{Solving Differential Equations in R.}
\end{itemize}

\subsection{Session info}\label{session-info}

\begin{Shaded}
\begin{Highlighting}[]

\end{Highlighting}
\end{Shaded}

\subsection{See also --- chapter index}\label{see-also-chapter-index}

\begin{itemize}
\tightlist
\item
  \href{../index.Rmd}{Methods in this chapter}
\item
  \href{../../../ch00-find-your-method.Rmd}{Find your method (Chapter
  0)}
\end{itemize}

\end{document}
